\documentclass[11pt,a4paper]{article}
\usepackage[margin=0.7in, top=0.6in, bottom=0.6in]{geometry}
\usepackage[utf8]{inputenc}
\usepackage[T1]{fontenc}
\usepackage{charter}
\usepackage{fontawesome5}
\usepackage{xcolor}
\usepackage{enumitem}
\usepackage{hyperref}
\usepackage{microtype}

\definecolor{accent}{HTML}{1E3A5F}     % Professional Dark Navy
\definecolor{accentblue}{HTML}{2B6CB0} % Slate Blue
\definecolor{lightgray}{HTML}{F3F4F6}
\definecolor{darktext}{HTML}{1F2937}
\definecolor{bulletline}{HTML}{9CA3AF}
\definecolor{subtletext}{HTML}{4B5563}

\hypersetup{colorlinks=true, linkcolor=accentblue, urlcolor=accentblue}

\setlength{\parindent}{0pt}
\setlength{\parskip}{8pt}

\begin{document}
\thispagestyle{empty}
\color{darktext}

% Encabezado idéntico al CV
\begin{center}
    {\Huge \bfseries\color{accent} SANTIAGO ALEJANDRO OCHOA QUESADA} \\[6pt]
    {\small\color{darktext}
    \faEnvelope\ \href{mailto:santy.alejo4@gmail.com}{santy.alejo4@gmail.com}
    \quad\textcolor{bulletline}{\textbar}\quad
    \faPhone\ +57 320 351 7153
    \quad\textcolor{bulletline}{\textbar}\quad
    \faMapMarker*\ Bogotá, Colombia \\[4pt]
    \faLinkedin\ \href{https://linkedin.com/in/santiago-alejandro-ochoa-quesada-a07a18288/}{LinkedIn}
    \quad\textcolor{bulletline}{\textbar}\quad
    \faGithub\ \href{https://github.com/Sochoaq}{GitHub}
    }
\end{center}
\vspace{2pt}
{\color{bulletline}\hrule height 0.8pt}
\vspace{10pt}

% Metadatos de la carta
\noindent
\begin{minipage}[t]{0.62\textwidth}
    \textbf{\color{accent} Para:}\\
    Equipo de Selección \& People\\
    \textbf{Mercado Libre Colombia}\\
    Bogotá, Colombia
\end{minipage}
\hfill
\begin{minipage}[t]{0.35\textwidth}\raggedleft
    \textbf{\color{accent} Fecha:}\\
    23 de octubre de 2026
\end{minipage}

\vspace{6pt}
\noindent
{\large\bfseries\color{accent} Asunto: Postulación al rol de Artificial Intelligence Engineer}
\vspace{2pt}

\noindent Estimado Equipo de Selección \& People,

Me dirijo a ustedes con gran entusiasmo para postularme al rol de \textbf{Artificial Intelligence Engineer} en Bogotá. Como Ingeniero Electrónico graduado de la \textbf{Universidad Nacional de Colombia} con un GPA de 4.6/5.0, poseo una sólida base matemática y científica, combinada con experiencia práctica en el desarrollo de software robusto, procesamiento de datos y algoritmos avanzados en Python.

A lo largo de mi trayectoria en AM4, he liderado el diseño y la implementación de sistemas inteligentes, automatizaciones industriales y arquitecturas de backend orientadas al análisis en tiempo real y la toma de decisiones basada en datos. Mi experiencia abarca desde el procesamiento de señales (DSP) y la visión artificial hasta la contenerización con Docker y la construcción de APIs escalables, garantizando siempre altos estándares de calidad, baja latencia y optimización de recursos.

Me motiva profundamente el propósito de Mercado Libre de democratizar el comercio electrónico y los servicios financieros en América Latina a través de la tecnología. Estoy convencido de que mi capacidad analítica, mi rapidez de aprendizaje y mi pasión por la ingeniería de inteligencia artificial me permitirán aportar valor inmediato en el diseño, experimentación y escalado de sistemas innovadores y seguros.

Agradezco de antemano su tiempo al considerar mi postulación y espero tener la oportunidad de conversar pronto sobre cómo mis competencias técnicas pueden contribuir al éxito de sus próximos desafíos tecnológicos.

\vspace{10pt}
\noindent Cordialmente,

\vspace{12pt}
\noindent
\textbf{\large\color{accent} Santiago Alejandro Ochoa Quesada}\\
{\small\color{subtletext} Ingeniero Electrónico | AI Engineering \& Software}\\
{\footnotesize\color{subtletext} santy.alejo4@gmail.com \;\textbar\; +57 320 351 7153 \;\textbar\; Bogotá, Colombia}

\end{document}
